\section{Métodos}
\label{sec:metodos}

\paragraph{Aumento de datos.}
En entrenamiento, cada recorte se rota hasta $\pm15^\circ$, se recorta al azar conservando entre
el 80 y el 100\,\% del área (\code{RandomResizedCrop} con aspecto 1:1), se voltea horizontalmente
con probabilidad 0,5 y se varían el brillo y el contraste ($\pm30\,\%$) y la saturación
($\pm20\,\%$) \citep{shorten2019augmentation}. Ninguna de estas transformaciones cambia el
significado del gesto. Por eso se excluye el volteo vertical, que convierte \cls{like} en
\cls{dislike}.

\paragraph{Arquitectura.}
GestureCNN (\cref{fig:arquitectura}) encadena cinco bloques. Cada uno aplica dos veces la
secuencia convolución $3\times3$, BN y ReLU, y termina con un \emph{max pooling} $2\times2$. Los
canales crecen de 32 a 256 mientras la resolución baja de $128^2$ a $4^2$. Los tres últimos
bloques añaden \code{Dropout2d} (0,1, 0,2 y 0,3), que apaga mapas de características completos
\citep{tompson2015spatialdropout}. La cabeza aplica \emph{global average pooling}, un
\emph{dropout} de 0,5 y una capa lineal de 256 a 9 salidas: 2\,313 parámetros, frente a los
36\,873 de una capa lineal sobre el mapa aplanado. Las variantes de la ablación quitan BN o el
\emph{dropout} y dejan intacto el resto de la red.

\begin{figure}[t]
  \centering
  \begin{tikzpicture}[
      font=\footnotesize,
      blk/.style={draw, rounded corners=2pt, minimum height=1.25cm, minimum width=1.45cm,
                  align=center, fill=#1},
      arr/.style={-{Stealth[length=2mm]}, thick}]
    \node[blk=gray!10]   (in) {Entrada\\$3\times128^2$};
    \node[blk=blue!10,   right=0.35cm of in]  (b1) {Bloque 1\\32 c.\\$\to 64^2$};
    \node[blk=blue!10,   right=0.35cm of b1]  (b2) {Bloque 2\\64 c.\\$\to 32^2$};
    \node[blk=blue!18,   right=0.35cm of b2]  (b3) {Bloque 3\\128 c.\\$\to 16^2$\\{\scriptsize drop 0,1}};
    \node[blk=blue!26,   right=0.35cm of b3]  (b4) {Bloque 4\\256 c.\\$\to 8^2$\\{\scriptsize drop 0,2}};
    \node[blk=blue!34,   right=0.35cm of b4]  (b5) {Bloque 5\\256 c.\\$\to 4^2$\\{\scriptsize drop 0,3}};
    \node[blk=orange!15, right=0.35cm of b5]  (gap) {GAP\\256};
    \node[blk=orange!15, right=0.35cm of gap] (fc) {Dropout 0,5\\Lineal\\$256\to9$};
    \foreach \a/\b in {in/b1,b1/b2,b2/b3,b3/b4,b4/b5,b5/gap,gap/fc} \draw[arr] (\a) -- (\b);
    \node[below=0.45cm of b3, draw, dashed, rounded corners=2pt, align=left, inner sep=4pt,
          font=\scriptsize] (det)
      {Cada bloque: $[\text{Conv }3{\times}3 \to \text{BN} \to \text{ReLU}]\times2 \to
       \text{MaxPool }2{\times}2\ (\to \text{Dropout2d})$};
    \draw[dashed, gray] (det.north) -- (b3.south);
    \node[above=0.12cm of b5, font=\scriptsize, text=red!70!black, align=center]
      {Grad-CAM: 2.ª conv. ($256\times8^2$)};
  \end{tikzpicture}
  \caption{Arquitectura de GestureCNN (\ParamsFinal{} parámetros): canales, resolución de salida
    y probabilidad de \code{Dropout2d} de cada bloque. Grad-CAM se calcula sobre la última
    convolución.}
  \label{fig:arquitectura}
\end{figure}

\paragraph{Entrenamiento.}
Usamos entropía cruzada sin pesos de clase, porque las clases están equilibradas, y Adam
\citep{kingma2015adam} con una tasa de aprendizaje inicial de $10^{-3}$, \emph{weight decay}
$10^{-4}$ y lotes de 64 (32 en el modelo final). La tasa se divide por dos tras 3 épocas sin mejorar
la pérdida de validación. La parada anticipada \citep{prechelt1998early} detiene el entrenamiento
tras 6 épocas sin mejora, con un máximo de 50, y restaura los pesos de la mejor época. La semilla
es 42 salvo que se indique otra. Todo se entrenó con PyTorch \VersionTorch{} en una \GPU{} de
portátil. Como referencia, ResNet18 preentrenado en ImageNet \citep{deng2009imagenet} se reajusta
entero con el mismo protocolo, una capa final de 9 salidas y una tasa diez veces menor
($10^{-4}$).

\subsection{Búsqueda de hiperparámetros}
\label{sec:busqueda-metodo}

Partiendo del modelo completo (tasa $10^{-3}$, lotes de 64, cinco bloques de 32 a 256 filtros,
sin capa oculta y ReLU), hicimos una búsqueda manual en la que cada configuración cambia
\emph{un único} hiperparámetro y mantiene fijo todo lo demás, incluidas la semilla y la partición.
Probamos \NConfigsBusqueda{} configuraciones:
\begin{itemize}
  \item tasa de aprendizaje de $3\cdot10^{-4}$ y de $3\cdot10^{-3}$;
  \item lotes de 32 y de 128;
  \item la mitad de filtros por bloque (\ParamsFiltrosMitad{} parámetros), cuatro bloques en lugar
    de cinco (\ParamsCuatroBloques) y una capa oculta de 128 neuronas en la cabeza;
  \item activaciones ELU \citep{clevert2016elu} y LeakyReLU (pendiente 0,1) en lugar de ReLU.
\end{itemize}
Cada configuración se entrenó una vez y se evaluó solo en validación. El criterio de elección,
fijado de antemano, fue la exactitud de validación. Cambiar un único factor cada vez permite
atribuirle la diferencia y cuesta una ejecución por valor probado, unos 20 minutos; a cambio, no
explora interacciones entre hiperparámetros.

\subsection{Protocolo de evaluación}
\label{sec:protocolo}

\paragraph{Selección del modelo.}
Todas las decisiones se tomaron con validación, pero el modelo final se eligió en dos fases. Lo
indicamos explícitamente porque afecta a la lectura del test. Primero, el modelo final fue la
variante completa, que era la mejor en validación de la ablación. Con ella, con las variantes y
con ResNet18 se evaluó el test una vez por modelo. Después hicimos la búsqueda de
hiperparámetros, solo con validación. La configuración con lotes de 32 fue la mejor y pasó a ser el
modelo final. La decisión se tomó antes de evaluar esa configuración en test, donde también se
evaluó una sola vez; pero el test ya se había mirado con el modelo anterior. Por eso presentamos
y comparamos los dos (\cref{sec:test}).

\paragraph{Métricas e incertidumbre.}
Reportamos la exactitud, la F1 macro y la matriz de confusión normalizada por filas, cuya
diagonal es el \emph{recall}. Como las imágenes de una misma persona no son independientes, los
intervalos de confianza del 95\,\% se calculan por \emph{bootstrap} remuestreando \emph{personas}
(\NBoot{} réplicas) \citep{efron1993bootstrap}. Para comparar dos modelos sobre las mismas
imágenes de test usamos la prueba exacta de McNemar \citep{mcnemar1947,dietterich1998tests}, y
repetimos la ablación con una segunda semilla (43) para medir la variación entre ejecuciones.
Estas comparaciones son exploratorias y no se corrigen por comparaciones múltiples. Grad-CAM
\citep{selvaraju2017gradcam} se calcula sobre la última convolución ($8\times8$) y siempre en
validación.
